\section*{Erd\H{o}s Problem 407}

\paragraph{Statement.}
Is the number \(w(n)\) of nonnegative integer quadruples satisfying
\[
 n=2^a+3^b+2^c3^d
\]
bounded by an absolute constant?

\paragraph{The external theorem and its counting convention.}
Bajpai and Bennett prove in Theorem 3 of
\emph{Effective \(S\)-unit equations beyond three terms: Newman's
conjecture} that the number \(\omega(n)\) of distinct sets
\(\{2^a,3^b,2^c3^d\}\) arising in such representations satisfies
\[
 \omega(n)\leq9\quad(n\geq1),\qquad
 \omega(n)\leq4\quad(n\geq131082).
 \tag{1}
\]
This established theorem is an explicit external dependency; its
Diophantine estimates and computational classification are not reproved
here. We check carefully that its convention answers the literal
quadruple-counting question.

\paragraph{Converting sets to quadruples.}
Fix a set \(S\) occurring in a representation of the fixed integer \(n\).
If \(|S|=3\), its three elements each occur once, so there are at most
six ordered triples with this underlying set.
If \(S=\{u,v\}\) with \(u\ne v\), the possible multisets are
\(\{u,u,v\}\) and \(\{u,v,v\}\). They have different sums, because
\(2u+v=u+2v\) would force \(u=v\). For our fixed \(n\), therefore,
at most one occurs, giving at most three ordered triples.
For \(|S|=1\) there is only one.

Each ordered triple \((x,y,z)\) of the required form determines
\((a,b,c,d)\) uniquely: \(x=2^a\), \(y=3^b\), and unique prime
factorization gives \(z=2^c3^d\). This remains true when one or more
entries are \(1\). Consequently
\[
 w(n)\leq6\omega(n)\leq54,
 \tag{2}
\]
and \(w(n)\leq24\) for \(n\geq131082\). This proves the requested
absolute boundedness, without incorrectly identifying the two counts.

\paragraph{Why a constant is the natural scale.}
For every \(a\geq5\), the integer \(n=2^a+9\) has representations with
the four distinct underlying sets
\[
 \{2^{a-1},9\},\quad
 \{2^{a-2},9,3\cdot2^{a-2}\},\quad
 \{2^a,3,6\},\quad
 \{2^a,1,8\}.
\]
Their respective sums, with \(2^{a-1}\) repeated in the first, are all
\(n\). The displayed sets are pairwise different for \(a\geq5\).
Thus infinitely many integers have at least four representations
even in the stricter set-counting convention.

\paragraph{Sources and scope.}
Primary theorem, including its definition of distinct representations:
\url{https://arxiv.org/abs/2308.05162}, Theorem 3.
Published version: Acta Arithmetica 214 (2024), 421--458,
\url{https://doi.org/10.4064/aa230725-14-9}.
Current statement: \url{https://www.erdosproblems.com/407},
checked 24 September 2026. This is a complete deduction using (1),
not an independent proof of that external theorem.

\paragraph{Completion estimate.}
The full boundedness claim follows from the cited theorem with the counting convention resolved.
\textbf{COMPLETION: 100\%}
